\begin{center}
\begin{tikzpicture}[x=1cm,y=1cm]
\draw[rounded corners=2pt,fill=blue!3,draw=blue!22] (0,0) rectangle (6.35,2.35);
\node[anchor=west,font=\small\bfseries,text=newsblue] at (.22,2.07) {甲烷中的氘／氢比：两颗彗星};
\node[anchor=east,font=\footnotesize] at (1.65,1.48) {3I/ATLAS};
\fill[newsblue!80] (1.82,1.34) rectangle (5.45,1.60);
\node[anchor=west,font=\footnotesize] at (5.51,1.47) {$3.18\%$};
\node[anchor=east,font=\footnotesize] at (1.65,.85) {67P};
\fill[orange!78] (1.82,.72) rectangle (2.10,.98);
\node[anchor=west,font=\footnotesize] at (2.17,.85) {$\approx0.245\%$};
\node[anchor=west,font=\scriptsize,text=newsgray] at (.24,.25) {67P 数值由 3I/ATLAS 比值及论文的 $13\pm2$ 倍关系推算。};
\end{tikzpicture}
\par\footnotesize\color{newsgray}图 1｜编辑部原创比例条。论文报告 3I/ATLAS 的甲烷 $\mathrm{D/H}=(3.18\pm0.30)\%$；67P 为近似反推值，非独立测量的原始报价。来源：Roth 等，\emph{Nature Astronomy}，2026-10-23。\color{black}
\end{center}